\section*{Chapter \ref{ch:nw:the-european-map} --- The European Map}

\begin{solution}{exo:nw:the-european-map:1}
From \qty{377.2}{\micro\second} (LD4 to Basildon, \qty{77.4}{\kilo\metre}) to \qty{4838.7}{\micro\second} (LD4 to Bergamo,
\qty{992.2}{\kilo\metre}): \qty{4.46}{\milli\second} more, one way.
\end{solution}

\begin{solution}{exo:nw:the-european-map:2}
One way, \qty{2.32}{\milli\second}; the floor in air over \qty{677.7}{\kilo\metre} is \qty{2.261}{\milli\second}; $2.320/2.261 = 1.026$,
\qty{59}{\micro\second} above the floor, and \qty{985}{\micro\second} under the straight-fibre floor.
\end{solution}

\begin{solution}{exo:nw:the-european-map:3}
\omterm{def:nw:colocation-products-and-how-they-are-sold:colo}{Colocation} puts the firm's equipment in the building where the venue matches, next to its engines; a \omterm{def:nw:the-european-map:pop}{point of presence} is an access point of
the venue's network in another building, from which the venue's network carries the firm's traffic to the engines. The first saves the
distance; the second only saves the firm from building its own long-distance link.
\end{solution}

\begin{solution}{exo:nw:the-european-map:4}
Antennas on a roof sit at different distances from the \omterm{def:nw:colocation-products-and-how-they-are-sold:mmr}{meet-me room}, and the cable from each antenna to the room would otherwise give the
nearest position an advantage of a few nanoseconds per metre (chapter~9). Equalising from the back of the antenna makes every wireless
provider start the building's race at the same point, which is what the equal-conditions rule of chapter~9 requires; without it, providers
would compete for roof positions rather than for straighter routes.
\end{solution}

\begin{solution}{exo:nw:the-european-map:5}
Through Slough: $678 + 992 = \qty{1670}{\kilo\metre}$, \qty{5.57}{\milli\second} in vacuum. Through Zurich: $307 + 205 = \qty{511}{\kilo\metre}$,
\qty{1.71}{\milli\second}. Direct: \qty{498}{\kilo\metre}, \qty{1.66}{\milli\second}. Zurich costs a \qty{13.6}{\kilo\metre} detour,
\qty{45}{\micro\second}; Slough more than triples the path.
\end{solution}

\begin{solution}{exo:nw:the-european-map:6}
It is an upper bound on a one-way time for orders from LD4 to IT3: the \omterm{def:nw:the-north-american-map:factor}{route factor} in air is at most $4/3.311 = 1.21$, at most
\qty{689}{\micro\second} above the floor, and at least \qty{839}{\micro\second} better than straight glass could be. It does not give the
typical figure, its distribution, the market-data direction, the behaviour in bad weather (the fibre back-up), or where exactly the clock
starts and stops.
\end{solution}

\begin{solution}{exo:nw:the-european-map:7}
Asking \texttt{primary} for XPAR on that day raises \texttt{LookupError}: the table gives two primary sites for that day. That is right: a venue matches in
one place, so two rows mean the table is wrong (an overlap in dates, a missing end date), and a silent choice would put a wrong number into
every floor computed from it.
\end{solution}

\begin{solution}{exo:nw:the-european-map:8}
The race that matters for a market maker on Euronext is among the firms quoting there, and any firm can colocate in Bergamo, including a
London firm. What the migration changed is the distance between Euronext and London's other venues: a firm quoting Euronext from Bergamo is
three milliseconds from its London hedges, and one quoting from London is three milliseconds from Euronext's book. It must choose where each
engine sits; it is not excluded.
\end{solution}

\begin{solution}{pb:nw:the-european-map:1}
\textbf{Part I.}
\begin{enumerate}
\item \qty{77.4}{\kilo\metre}: \qty{258.0}{\micro\second} in vacuum, \qty{377.2}{\micro\second} in straight fibre.
\item \qty{603.2}{\kilo\metre}: \qty{2012}{\micro\second} and \qty{2942}{\micro\second}.
\item A sliver: 678, 603 and \qty{77}{\kilo\metre}; Euronext was almost in London.
\item Anywhere in the London triangle, LD4 among them: a few hundred microseconds from Basildon, the Docklands and the London points of
  presence.
\end{enumerate}
\textbf{Part II.}
\begin{enumerate}[resume]
\item \qty{992.2}{\kilo\metre}: \qty{3309.7}{\micro\second} in vacuum, \qty{4838.7}{\micro\second} in fibre.
\item $+\qty{3052}{\micro\second}$ in vacuum and $+\qty{4461}{\micro\second}$ in fibre.
\item \qty{497.6}{\kilo\metre}, \qty{1659.9}{\micro\second} in vacuum: \qty{352}{\micro\second} less than from Basildon
  (\qty{515}{\micro\second} less in fibre).
\item Wide: 678, 498 and \qty{992}{\kilo\metre}; Frankfurt is now nearer to Euronext than London is.
\end{enumerate}
\textbf{Part III.}
\begin{enumerate}[resume]
\item $1.2 \times 4\,838.7 = \qty{5806}{\micro\second}$.
\item Less than \qty{4000}{\micro\second}: at least \qty{1806}{\micro\second} saved.
\item For quoting Euronext's book, yes: its cancels and requotes then travel metres, not a thousand kilometres, and Eurex is nearer to Bergamo
  (\qty{1.66}{\milli\second}) than to LD4 (\qty{2.26}{\milli\second}). It loses nearness to London's venues (about \qty{3.1}{\milli\second}
  from Bergamo) and its London staff's easy access; many firms run both.
\item $306.7 + 204.5 - 497.6 = \qty{13.6}{\kilo\metre}$, \qty{45}{\micro\second} of light in vacuum.
\end{enumerate}
\textbf{Part IV.}
\begin{enumerate}[resume]
\item \emph{Named result}: the one-way floor from Slough LD4 to Euronext's matching engines rose from \qty{0.26}{\milli\second} to
  \qty{3.31}{\milli\second} in vacuum (\qty{+3.05}{\milli\second}) and from 0.38 to \qty{4.84}{\milli\second} in fibre
  (\qty{+4.46}{\milli\second}); the London--Frankfurt--Euronext triangle went from 678, 603 and \qty{77}{\kilo\metre} to 678, 498 and
  \qty{992}{\kilo\metre}.
\item So that London clients could connect to Bergamo, test and move before the cutover, without building their own links to Italy.
\item Two rows, Basildon until 5 June 2022 and Bergamo from 6 June, each with its source: a query by date then gives the right answer for any
  day, and history is kept for backtests that must use the latencies of their own dates.
\item The London Stock Exchange's move from the City to the Docklands, planned for October 2022 and made in February 2023.
\item That a relay should lie on or near the Frankfurt--Bergamo geodesic; Zurich costs \qty{45}{\micro\second}, London \qty{3.9}{\milli\second}.
\item Basildon, Bergamo and Upplands Väsby. An error of a kilometre or two moves a floor by a few microseconds, against a change of three
  thousand: not for this result.
\item A fibre whose light travels nearly as in air would bring the fibre floor close to the air floor, \qty{3.3}{\milli\second}, removing most of
  microwave's advantage on this corridor, with fibre's bandwidth and weather-independence (chapter~13).
\item It moved Euronext three milliseconds away and made the firm choose, engine by engine, between being near Euronext and being near London.
\end{enumerate}
\end{solution}

\begin{solution}{iq:nw:the-european-map:1}
Euronext in Bergamo (Aruba IT3), Eurex and Xetra in Frankfurt (Equinix FR2), the London Stock Exchange in the Docklands, ICE Futures Europe in
Basildon, SIX in Zurich (ZH4), Nasdaq Nordic near Stockholm. London (Slough) to Frankfurt is about \qty{678}{\kilo\metre}: \qty{2.26}{\milli\second}
one way in air, \qty{3.3}{\milli\second} in straight fibre.
\iqlookfor{Buildings rather than cities; the 2022 Euronext move; one-way floors in both media.}
\end{solution}

\begin{solution}{iq:nw:the-european-map:2}
An access point of the venue's network in a building where its clients already are, so that they connect with a \omterm{def:nw:colocation-products-and-how-they-are-sold:mmr}{cross-connect} and the venue
carries the traffic to its engines. An exchange opens one abroad when its clients are there: Euronext opened two in London before moving to
Italy.
\iqlookfor{PoP against \omterm{def:nw:colocation-products-and-how-they-are-sold:colo}{colocation}; the venue's network and its latency; resilience (two PoPs); the \omterm{def:nw:colocation-products-and-how-they-are-sold:mmr}{cross-connect}.}
\end{solution}

\begin{solution}{iq:nw:the-european-map:3}
Keep it as data with a source and a read date per row, one row per venue, role and period; re-verify against the venues' notices on a
schedule; flag rows older than a limit; never overwrite a row when a venue moves, close it and add another.
\iqlookfor{Provenance; history kept for backtests; automated staleness checks; subscription to venue notices.}
\end{solution}

\begin{solution}{iq:nw:the-european-map:4}
Read the venue's migration documents; compute the new floors from each of the firm's sites; decide where each strategy's engine should sit;
order space, power and connectivity in the new building and links to it, with lead times; test in the venue's rehearsals; plan the cutover
weekend and the fallback; update the site registry and every latency assumption in research.
\iqlookfor{Order of operations and lead times; the rehearsal; research assumptions updated, not just production.}
\end{solution}

\begin{solution}{iq:nw:the-european-map:5}
Because their clients are elsewhere and their own building's appeal depends on how fast those clients can reach it; a venue-run microwave
link with equal access (SIX from Zurich to Frankfurt, Euronext from London to Bergamo) sells speed under the venue's fair-access rules and
brings order flow.
\iqlookfor{Fair access; revenue; competition with third-party networks; one direction or both.}
\end{solution}

\begin{solution}{iq:nw:the-european-map:6}
The floors say: near the Frankfurt--Bergamo line, or in one of the two buildings with a fast link to the other; Zurich is only
\qty{45}{\micro\second} off the direct line. Check with the registry and the geodesics, then with published routes and their \omterm{def:nw:the-north-american-map:factor}{route factors},
and finally by measuring a candidate link with timestamped captures at both ends.
\iqlookfor{Triangle inequality; which side's latency matters more for the strategy; measurement, not only floors.}
\end{solution}
